\section{Related Work}
The standard approach for addressing non-stationarity in MARL is to consider information about the other agents and reason about the effects of their joint actions~\citep{hernandezLealK17survey}.
The literature on opponent modeling, for instance, infers opponents' behaviors and conditions an agent's policy on the inferred behaviors of others~\citep{he16opponent-modeling,raileanu18opponent-modeling,grover18policy-representation}.
Studies regarding the centralized training with decentralized execution framework~\citep{lowe17maddpg,foerster2017counterfactual,yang18mean-field-marl,wen2018probabilistic}, which accounts for the behaviors of others through a centralized critic, can also be classified into this category.
While this body of work alleviates non-stationarity, it is generally assumed that each agent will have a stationary policy in the future.
Because other agents can have different behaviors in the future as a result of learning~\citep{foerster17lola}, this incorrect assumption can cause sample inefficient and improper adaptation.
In contrast, Meta-MAPG models each agent's learning process, allowing a meta-learning agent to adapt efficiently.

Our approach is also related to prior work that considers the learning of other agents in the environment. 
This includes~\citet{zhang10lookahead} who attempted to discover the best response adaptation to the anticipated future policy of other agents. Our work is also related, as discussed previously, to LOLA \citep{foerster17lola} and more recent improvements~\citep{foerster2018dice}.
Another relevant idea explored by~\citet{letcher2018stable} is to interpolate between the frameworks of~\citet{zhang10lookahead} and~\citet{foerster17lola} in a way that guarantees convergence while influencing the opponent's future policy.
However, all of these approaches only account for the learning processes of other agents and fail to consider an agent's own non-stationary policy dynamics as in the own learning gradient discussed in the previous section. 
Additionally, these papers do not leverage meta-learning. As a result, these approaches may require many samples to properly adapt to new agents.

Meta-learning \citep{schmidhuber87meta_learning, bengio1992optimization} has recently become very popular as a method for improving sample efficiency in the presence of changing tasks in the deep RL literature \citep{wang2016learning,duan2016rl,maml,snail,Reptile}. See~\citet{vilalta2002perspective} and~\citet{hospedales2020meta} for in-depth surveys of meta-learning. In particular, our work builds on the popular model-agnostic meta-learning framework~\citep{maml}, where gradient-based learning is used both for conducting so called inner-loop learning and to improve this learning by computing gradients through the computational graph. When we train our agents so that the inner loop can accommodate for a dynamic Markov chain of other agent policies, we are leveraging an approach that has recently become popular for supervised learning called meta-continual learning~\citep{MER,Javed2019Meta,spigler2019meta,beaulieu2020learning,caccia2020online,Gupta2020LaMAMLLM}. 
This means that our agent trains not just to adapt to a single set of policies during meta-training, but rather to adapt to a set of changing policies with Markovian updates. 
As a result, we avoid an issue of past work \citep{alshedivat2018continuous} that required the use of importance sampling during meta-testing (see~\Cref{sec:meta-pg-difference-details} for more details).